% ##############################################################################
\section{Toward a decentralized, anonymized \NoesisMarket{}}
\label{sec:decentralized_extension}

Sections~\ref{sec:noesis_market} and \ref{sec:noesis_server} describe \Noesis{}
as operated by a single, trusted market and gateway. That design is
deliberately the starting point: it lets the auction
(Section~\ref{sec:noesis_market}) and the fulfillment layer
(Section~\ref{sec:noesis_server}) be specified without also solving
distributed-systems and cryptographic-protocol problems.
%
This section sketches an alternative in which \NoesisMarket's auction runs on a
public blockchain rather than inside a single operator's infrastructure, with
bids and asks kept anonymous until a round closes. The design adapts ideas
from stake-based, blind-bid marketplaces for off-chain computation, in
which a requester who needs a unit of work performed under a deadline is
matched against a supplier who is paid for performing it:
\NoesisMarket's clearing problem (Problem~\ref{prob:clearing}) shares the
same shape as that matching problem, so the mechanisms below carry over
with only minor changes. Section~\ref{sec:sota_decentralized} surveys
production systems, Bittensor, Gensyn, and Ritual among them, that already
run parts of this machinery, validator scoring, proof-of-learning,
proof-of-execution, at network scale; the question this section asks is
how far analogous machinery can go when applied to \NoesisMarket's own
auction and contract, rather than to a network-internal reward schedule.

We treat what follows as a design sketch rather than a specification.
Section~\ref{sec:decentralized_open_questions} collects the open
questions it leaves unresolved.

The interface \Noesis{} exposes to buyers and sellers, submitting a bid
or an ask and receiving a contract (Section~\ref{sec:protocol_overview}),
does not need to change under this extension: what changes is where and
how a round is cleared and settled, not the shape of what a participant
submits or receives.

% ==============================================================================
\subsection{Motivation: censorship resistance and bidder privacy}
\label{sec:decentralized_motivation}

The centralized design of Section~\ref{sec:noesis_market} concentrates two
kinds of trust in a single operator.

\begin{itemize}
  \item \textbf{Visibility}: the operator sees every bid and ask,
    including a buyer's required capability tier and volume, which can
    reveal competitively sensitive information, for instance that a
    buyer is about to run a large batch job, before the round even
    clears.
  \item \textbf{Discretion}: the operator alone decides which asks are
    eligible to compete, which contracts clear, and which sellers are
    excluded on reputation grounds (Section~\ref{sec:reputation}), with
    no mechanism for a participant to independently verify that this was
    done correctly.
\end{itemize}

A blockchain-based \NoesisMarket{} addresses both concerns by moving the
parts of the mechanism that most need to be trust-minimized, contract
terms, matching, settlement, and non-performance penalties, onto a public
ledger, while keeping bid and ask contents hidden until the round closes
(Section~\ref{sec:blind_bidding}). Neither property is free: on-chain
settlement trades the negligible latency and cost of a centralized
clearinghouse for the latency, gas cost, and finality assumptions of the
underlying chain, a trade-off Section~\ref{sec:decentralized_comparison}
makes explicit.

Figure~\ref{fig:trust_boundaries} contrasts the two trust boundaries
before the mechanisms below are introduced: in the centralized design, a
single operator sees every bid and ask in full and alone decides
eligibility and matching, while in the decentralized design, only
on-chain addresses are visible until the reveal window closes
(Section~\ref{sec:blind_bidding}), and matching and eligibility are
verifiable on-chain rather than taken on the operator's word.

% rendered_images:begin
% ```graphviz
% digraph TrustBoundaries {
%   // ---------------------------------------------------------------
%   bgcolor="transparent";
%   pad="0.08";
%   splines=spline;
%   nodesep=0.35;
%   ranksep=0.42;
%   rankdir=TB;
% 
%   node [shape=box,
%         style="rounded,filled",
%         fillcolor="#FFFFFF",
%         color="#9AA9B8",
%         penwidth=1.8,
%         fontname="Helvetica",
%         fontcolor="#243B53",
%         fontsize=10,
%         margin="0.20,0.12",
%         height=0.42,
%         width=2.4];
% 
%   edge [color="#A3B1C0",
%         penwidth=1.3,
%         arrowhead=vee,
%         arrowsize=0.75,
%         fontname="Helvetica",
%         fontsize=8.5,
%         fontcolor="#7B8794"];
% 
%   // ---------------------------------------------------------------
%   // The two clusters are stacked vertically (via the invisible
%   // connecting edge below) to keep the figure narrow.
%   // source : warm sand   |   actor : rose   |   external : violet   |   process : blue
% 
%   subgraph cluster_centralized {
%     label     = "Centralized";
%     labelloc  = "t";
%     fontname  = "Helvetica-Bold";
%     fontsize  = 10.5;
%     fontcolor = "#3A4A5C";
%     style     = "rounded,filled";
%     fillcolor = "#F7F9FB";
%     color     = "#C7D0DA";
%     penwidth  = 1.0;
%     margin    = 12;
% 
%     BidsC    [label=<<b>Full bid/ask contents</b><br/><font point-size="8.5" color="#946A2E">Tier, volume, price</font>>,
%               fillcolor="#FBEBD4", color="#D9A85F", fontcolor="#6B4517"];
%     Operator [label=<<b>Single operator</b>>,
%               fillcolor="#F6E1E8", color="#D98CA8", fontcolor="#6B2A44"];
%     Decide   [label=<<b>Decides eligibility</b><br/><font point-size="8.5" color="#41719C">&amp; matching</font>>,
%               fillcolor="#D3E3F3", color="#7CA6CE", fontcolor="#1F4E79"];
%     BidsC -> Operator [label="Sees all", penwidth=1.6, color="#8592A3"];
%     Operator -> Decide [penwidth=1.6, color="#8592A3"];
%   }
% 
%   subgraph cluster_decentralized {
%     label     = "Decentralized";
%     labelloc  = "t";
%     fontname  = "Helvetica-Bold";
%     fontsize  = 10.5;
%     fontcolor = "#3A4A5C";
%     style     = "rounded,filled";
%     fillcolor = "#F7F9FB";
%     color     = "#C7D0DA";
%     penwidth  = 1.0;
%     margin    = 12;
% 
%     BidsD  [label=<<b>Only addresses visible</b><br/><font point-size="8.5" color="#946A2E">until reveal window</font>>,
%             fillcolor="#FBEBD4", color="#D9A85F", fontcolor="#6B4517"];
%     Chain  [label=<<b>Public chain</b><br/><font point-size="8.5" color="#6B52A8">Commit-reveal</font>>,
%             fillcolor="#E9E1F7", color="#A88FD9", fontcolor="#45296B"];
%     Verify [label=<<b>Matching &amp; eligibility</b><br/><font point-size="8.5" color="#41719C">verifiable on-chain</font>>,
%             fillcolor="#D3E3F3", color="#7CA6CE", fontcolor="#1F4E79"];
%     BidsD -> Chain [label="Commitments", penwidth=1.6, color="#8592A3"];
%     Chain -> Verify [penwidth=1.6, color="#8592A3"];
%   }
% 
%   Decide -> BidsD [style=invis];
% }
% ```
% label=fig:trust_boundaries
% caption=The two trust boundaries: centralized versus decentralized \NoesisMarket{}.
% rendered_images:end
% render_images:begin
\begin{figure}[H]
  \includegraphics[width=\linewidth]{08_decentralized_extension.tex.figs/08_decentralized_extension.1.png}
  \caption{The two trust boundaries: centralized versus decentralized \NoesisMarket{}.}
  \label{fig:trust_boundaries}
\end{figure}
% render_images:end

% ==============================================================================
\subsection{On-chain contracts, blind bidding, and anonymity}
\label{sec:blind_bidding}

The contract, bid, and ask notation of Section~\ref{sec:contracts_and_notation}
carries over unchanged: a decentralized round still clears
$(\ntasksattr, \clevelattr, \lmaxattr, \rminattr, \priceattr)$ tuples
(Definition~\ref{def:contract}). What changes is how a bid or an ask is
submitted and revealed.

\begin{definition}[Blind bid]
\label{def:blind_bid}
A \emph{blind bid} for a bid $\beta = (n_\beta, c_\beta, \ell_\beta,
r_\beta, p_\beta)$ (Definition~\ref{def:bid}) is a commitment
$$
  h_\beta = H(\beta \,\|\, \nu_\beta)
$$
posted on-chain during a round's \emph{commit window}, where $H$ is a
collision-resistant hash function and $\nu_\beta$ is a secret salt known
only to the bidder. During the round's subsequent \emph{reveal window},
the bidder discloses $(\beta, \nu_\beta)$; only bids and asks revealed
within that window, and matching a previously posted commitment, are
eligible for the clearing problem (Problem~\ref{prob:clearing}) run at the
round's close. An ask $\alpha$ is committed and revealed symmetrically.
\end{definition}

Definition~\ref{def:blind_bid} follows the standard commit-reveal pattern
for blind bidding, motivated by the same concern that motivates frequent
batch auctions over continuous limit order books
(Section~\ref{sec:related_auctions}): a party that can see pending
bids before a round closes can front-run them, either by submitting a
marginally better ask to snipe a favorable price, or, at the
block-production layer, by reordering or censoring specific bids.
Committing before revealing reduces the incentive to do so, at the cost of
one additional round-trip and the mempool-monitoring risk flagged in
Section~\ref{sec:decentralized_open_questions}: a party can still observe
bids as they are revealed and submit a competing bid before the reveal
window closes, unless that window is itself no longer than a single
block.

Anonymity is a byproduct of the same design: a buyer and a seller are
identified on-chain only by an address, and, until the reveal window, not
even the contents of their bid are visible. This does not, by itself,
resolve \NoesisServer's own privacy question, what safeguards are needed
before logging real prompt/response pairs (Section~\ref{sec:gateway_architecture}):
fulfillment still requires a buyer's requests to reach a real seller
off-chain, so anonymity here covers the auction, not the inference traffic
that a cleared contract subsequently generates.

Figure~\ref{fig:commit_reveal_sequence} traces one round through these
three phases: during the commit window, bidders post the commitment
$h_\beta = H(\beta \,\|\, \nu_\beta)$ (Definition~\ref{def:blind_bid});
during the reveal window, they disclose $(\beta, \nu_\beta)$; and only
revealed bids and asks matching a posted commitment are kept for the
clearing problem (Problem~\ref{prob:clearing}) run at the round's close.

% rendered_images:begin
% ```graphviz
% digraph CommitReveal {
%   // ---------------------------------------------------------------
%   bgcolor="transparent";
%   pad="0.06";
%   splines=spline;
%   nodesep=0.35;
%   ranksep=0.50;
% 
%   node [shape=box,
%         style="rounded,filled",
%         fillcolor="#FFFFFF",
%         color="#9AA9B8",
%         penwidth=1.8,
%         fontname="Helvetica",
%         fontcolor="#243B53",
%         fontsize=10,
%         margin="0.22,0.14",
%         height=0.50];
% 
%   edge [color="#A3B1C0",
%         penwidth=1.3,
%         arrowhead=vee,
%         arrowsize=0.75,
%         fontname="Helvetica",
%         fontsize=8.5,
%         fontcolor="#7B8794"];
% 
%   // ---------------------------------------------------------------
%   rankdir=TB;
% 
%   // source : warm sand   |   process : blue   |   artifact : sage green
% 
%   Commit [label=<<b>Commit window</b><br/><font point-size="8.5" color="#946A2E">Bidders post h = H(bid || salt)</font>>,
%           fillcolor="#FBEBD4", color="#D9A85F", fontcolor="#6B4517"];
%   Reveal [label=<<b>Reveal window</b><br/><font point-size="8.5" color="#41719C">Bidders disclose (bid, salt)</font>>,
%           fillcolor="#D3E3F3", color="#7CA6CE", fontcolor="#1F4E79"];
%   Match  [label=<<b>Commitment check</b><br/><font point-size="8.5" color="#4F7A5A">Only matching, revealed bids kept</font>>,
%           fillcolor="#DFEDE0", color="#8FB79A", fontcolor="#2E5A3D"];
%   Clear  [label=<<b>Clearing</b><br/><font point-size="8.5" color="#41719C">NoesisMarket clearing problem</font>>,
%           fillcolor="#D3E3F3", color="#7CA6CE", fontcolor="#1F4E79"];
% 
%   Commit -> Reveal -> Match -> Clear [penwidth=1.6, color="#8592A3"];
% }
% ```
% label=fig:commit_reveal_sequence
% caption=One auction round under commit-reveal.
% rendered_images:end
% render_images:begin
\begin{figure}[H]
  \includegraphics[width=\linewidth]{08_decentralized_extension.tex.figs/08_decentralized_extension.2.png}
  \caption{One auction round under commit-reveal.}
  \label{fig:commit_reveal_sequence}
\end{figure}
% render_images:end

% ==============================================================================
\subsection{Stake-backed fulfillment: slashing as decentralized reputation}
\label{sec:stake_slashing}

Section~\ref{sec:reputation}'s reputation score $\rho_\alpha(t)$ is a soft,
ex post penalty enforced by a trusted operator: a seller who under-delivers
is downgraded, but only the operator can act on that downgrade. A
decentralized \NoesisMarket{} replaces, or complements, this with a hard,
cryptoeconomic penalty.

\begin{definition}[Staked ask]
\label{def:staked_ask}
A \emph{staked ask} extends an ask $\alpha$ (Definition~\ref{def:ask})
with a stake $\sigma_\alpha \in \bbR_+$, denominated in a fixed token
$\tau$, posted to an escrow contract at the time $\alpha$ is revealed. If
$\alpha$ is matched into a contract $\kappa$, the stake is locked until
the fulfillment window $W$ closes (Section~\ref{sec:fulfillment_monitoring}).
It is released in full if $\kappa$ is not flagged as violated, i.e.,
$\hat r_{\text{lower}}(\kappa, W) \geq \rminattr(\kappa)$
(Equation~\eqref{eq:reliability_lower_bound}), and slashed, in full or in
part, otherwise.
\end{definition}

Definition~\ref{def:staked_ask} still needs the report of
$\hat r_{\text{lower}}(\kappa, W)$ itself, computed off-chain by
\NoesisServer{} from delivered responses; slashing does not remove the
need for fulfillment monitoring, it changes what happens once a violation
is detected, from a reputation downgrade to an immediate, automatic,
on-chain forfeiture. This is a stronger deterrent against the capability
misrepresentation risk of Section~\ref{sec:mechanism_design_risks}
precisely because it does not require the seller's continued
participation in future rounds to take effect. How $\hat
r_{\text{lower}}(\kappa, W)$ is reported on-chain in a form a smart
contract can act on, rather than treating \NoesisServer{} itself as a
trusted oracle, is left open in Section~\ref{sec:decentralized_open_questions}.

Figure~\ref{fig:staked_ask_lifecycle} shows this lifecycle as a state
diagram: the stake is posted when $\alpha$ is revealed, locked once
matched into a contract $\kappa$, and released in full or slashed once
$\hat r_{\text{lower}}(\kappa, W)$ is reported against $\rminattr(\kappa)$
at the close of the fulfillment window.

% rendered_images:begin
% ```graphviz
% digraph StakedAskLifecycle {
%   // ---------------------------------------------------------------
%   bgcolor="transparent";
%   pad="0.06";
%   splines=spline;
%   nodesep=0.35;
%   ranksep=0.70;
% 
%   node [shape=box,
%         style="rounded,filled",
%         fillcolor="#FFFFFF",
%         color="#9AA9B8",
%         penwidth=1.8,
%         fontname="Helvetica",
%         fontcolor="#243B53",
%         fontsize=10,
%         margin="0.22,0.14",
%         height=0.50];
% 
%   edge [color="#A3B1C0",
%         penwidth=1.3,
%         arrowhead=vee,
%         arrowsize=0.75,
%         fontname="Helvetica",
%         fontsize=8.5,
%         fontcolor="#7B8794"];
% 
%   // ---------------------------------------------------------------
%   rankdir=LR;
% 
%   // source : warm sand   |   artifact : sage green   |   risk : red
% 
%   Posted   [label=<<b>Posted</b><br/><font point-size="8.5" color="#946A2E">At reveal</font>>,
%             fillcolor="#FBEBD4", color="#D9A85F", fontcolor="#6B4517"];
%   Locked   [label=<<b>Locked</b><br/><font point-size="8.5" color="#4F7A5A">Matched into contract κ</font>>,
%             fillcolor="#DFEDE0", color="#8FB79A", fontcolor="#2E5A3D"];
%   Released [label=<<b>Released in full</b><br/><font point-size="8.5" color="#4F7A5A">Compliant</font>>,
%             fillcolor="#DFEDE0", color="#8FB79A", fontcolor="#2E5A3D"];
%   Slashed  [label=<<b>Slashed, in full or in part</b><br/><font point-size="8.5" color="#A83D3D">Violated</font>>,
%             fillcolor="#F8DCDC", color="#D97A7A", fontcolor="#6B1F1F"];
% 
%   Posted -> Locked [label="Matched", penwidth=1.6, color="#8592A3"];
%   Locked -> Released [label="Window W closes,\nr_lower >= R_min"];
%   Locked -> Slashed [
%     label="Window W closes,\nr_lower < R_min",
%     color="#C0455B", fontcolor="#C0455B"
%   ];
% }
% ```
% label=fig:staked_ask_lifecycle
% caption=Lifecycle of a staked ask.
% rendered_images:end
% render_images:begin
\begin{figure}[H]
  \includegraphics[width=\linewidth]{08_decentralized_extension.tex.figs/08_decentralized_extension.3.png}
  \caption{Lifecycle of a staked ask.}
  \label{fig:staked_ask_lifecycle}
\end{figure}
% render_images:end

Sizing $\sigma_\alpha$ raises a familiar asymmetry in per-bid staking: a
buyer setting the required stake unilaterally can grief sellers by
demanding collateral far in excess of
the contract's value, while a buyer who is not itself bonded has nothing
at risk if a matched contract turns out not to be genuine.
Section~\ref{sec:decentralized_open_questions} returns to this.

% ==============================================================================
\subsection{Where the clearing problem runs}
\label{sec:decentralized_matching_engine}

Problem~\ref{prob:clearing} is a linear program over the revealed bids and
asks in a tier bucket. Solving it naively on-chain, recomputing the
clearing price and every matched quantity $x_{\beta\alpha}$ inside a smart
contract, scales the gas cost of a round with $|B_c| \cdot |A_c|$, which
is unattractive for a bucket with more than a handful of participants.
Three placements are possible.

\begin{itemize}
  \item \textbf{Fully on-chain}: the clearing problem is solved inside the
    smart contract itself. Trust-minimized by construction, but the gas
    cost above makes it impractical beyond small buckets.

  \item \textbf{Fully off-chain}: a designated operator solves
    Problem~\ref{prob:clearing} off-chain and posts only the result,
    $p^*(c, t)$ and the matched contracts. Cheap, but reintroduces the
    single trusted operator that Section~\ref{sec:decentralized_motivation}
    set out to remove: nothing on-chain prevents the operator from posting
    an incorrect or self-favoring solution.

  \item \textbf{Verifiably off-chain}: an off-chain solver computes
    $p^*(c, t)$ and $\{x_{\beta\alpha}\}$ as before, but additionally
    produces a succinct proof $\phi$ that the solution satisfies the
    feasibility constraints of Problem~\ref{prob:clearing}
    (Equations~\eqref{eq:buyer_quantity}--\eqref{eq:price_limit}) and
    leaves no compatible pair unmatched at a price both sides would
    accept. A smart contract verifies $\phi$ in time independent of
    $|B_c| + |A_c|$ and only then finalizes the round.
\end{itemize}

The third option, running the matchmaking engine partially off-chain with
a zero-knowledge property showing the solution is correct, is the option
that preserves both the trust-minimization goal of
Section~\ref{sec:decentralized_motivation} and the practicality of the
fully-off-chain design. It is also the least developed of the three:
producing a succinct, efficiently verifiable proof of optimality for a
linear assignment problem, rather than merely of feasibility, is itself an
open implementation question, listed in
Section~\ref{sec:decentralized_open_questions}.

Figure~\ref{fig:clearing_placement_spectrum} spans the three placements
along this spectrum, from trust-minimized and expensive to cheap and
trusted: the verifiably-off-chain placement is the least developed of the
three but is the only one that preserves trust-minimization without the
gas cost of solving Problem~\ref{prob:clearing} on-chain.

% rendered_images:begin
% ```tikz
% \definecolor{axiscolor}{RGB}{133,146,163}
% \definecolor{ptcolor}{RGB}{31,78,121}
% \fontfamily{phv}\selectfont
% \draw[very thick, ->, >=stealth, axiscolor] (0.3,0) -- (10.7,0);
% \node[anchor=west, align=left, font=\footnotesize] at (0,-0.9) {trust-minimized,\\expensive};
% \node[anchor=east, align=right, font=\footnotesize] at (11,-0.9) {cheap,\\trusted};
% \filldraw[ptcolor] (1.2,0) circle (3pt);
% \node[align=center, font=\footnotesize] at (1.2,1.1) {fully on-chain\\(gas cost scales\\with $|B_c|\cdot|A_c|$)};
% \filldraw[ptcolor] (5.5,0) circle (3pt);
% \node[align=center, font=\footnotesize] at (5.5,1.4) {verifiably off-chain\\(succinct proof of\\feasibility; least developed)};
% \filldraw[ptcolor] (9.8,0) circle (3pt);
% \node[align=center, font=\footnotesize] at (9.8,1.1) {fully off-chain\\(operator posts\\the result only)};
% ```
% label=fig:clearing_placement_spectrum
% caption=The three placements for the clearing problem, on a trust/cost spectrum.
% rendered_images:end
% render_images:begin
\begin{figure}[H]
  \includegraphics[width=\linewidth]{08_decentralized_extension.tex.figs/08_decentralized_extension.4.png}
  \caption{The three placements for the clearing problem, on a trust/cost spectrum.}
  \label{fig:clearing_placement_spectrum}
\end{figure}
% render_images:end

% ==============================================================================
\subsection{Redundancy as a decentralization safeguard}
\label{sec:decentralized_redundancy}

A single seller's downtime is a single point of failure for the contract
that seller was matched into, the same risk a redundancy parameter
addresses in proof-generation marketplaces where a single prover might
fail to deliver a proof in time: a requester can choose to have
$k_{\text{red}} \geq 1$ provers race for the same job, each staking
collateral, with only the first, or, in an auction variant, the top few,
to deliver being paid.

The same parameter adapts directly to a decentralized \NoesisMarket: a
contract can be matched against $k_{\text{red}}$ staked asks instead of
one, with traffic split or raced across them during the fulfillment
window, and only the sellers who meet the contract's terms retaining their
stake. This is a different lever from answer fusion
(Section~\ref{sec:fusion}): fusion queries several providers per request
to improve the quality of a single answer, while redundancy here queries
several providers to hedge against any one of them being unavailable or
slashed, and is a settlement-layer property rather than an answer-quality
one. Redundancy also gives the market a tool against the concentration
risk noted in Section~\ref{sec:mechanism_design_risks}'s discussion of
thin, collusion-prone tier buckets: instead of always matching a buyer's
full volume against the single lowest-priced compatible ask, a
decentralized \NoesisMarket{} can split it across the $k_{\text{red}}$
lowest-priced compatible asks, so that the same top seller does not
deterministically win every round in a thin bucket.

Figure~\ref{fig:redundancy_racing} contrasts a redundant match with the
single-ask baseline: instead of matching a contract against a single
staked ask, it is matched against $k_{\text{red}}$ staked asks that race
to fulfill it, and only the sellers who meet the contract's terms retain
their stake, hedging against any one seller being unavailable or slashed.

% rendered_images:begin
% ```graphviz
% digraph Redundancy {
%   // ---------------------------------------------------------------
%   bgcolor="transparent";
%   pad="0.06";
%   splines=spline;
%   nodesep=0.35;
%   ranksep=0.50;
% 
%   node [shape=box,
%         style="rounded,filled",
%         fillcolor="#FFFFFF",
%         color="#9AA9B8",
%         penwidth=1.8,
%         fontname="Helvetica",
%         fontcolor="#243B53",
%         fontsize=10,
%         margin="0.22,0.14",
%         height=0.50];
% 
%   edge [color="#A3B1C0",
%         penwidth=1.3,
%         arrowhead=vee,
%         arrowsize=0.75,
%         fontname="Helvetica",
%         fontsize=8.5,
%         fontcolor="#7B8794"];
% 
%   // ---------------------------------------------------------------
%   rankdir=TB;
% 
%   // artifact : sage green   |   external (staked ask) : violet
% 
%   subgraph cluster_single {
%     label     = "Single-ask match (k_red = 1)";
%     labelloc  = "t";
%     fontname  = "Helvetica-Bold";
%     fontsize  = 10.5;
%     fontcolor = "#3A4A5C";
%     style     = "rounded,filled";
%     fillcolor = "#F7F9FB";
%     color     = "#C7D0DA";
%     penwidth  = 1.0;
%     margin    = 14;
% 
%     ContractS [label=<<b>contract κ</b>>,
%                fillcolor="#DFEDE0", color="#8FB79A", fontcolor="#2E5A3D"];
%     AskS      [label=<<b>one staked ask</b>>,
%                fillcolor="#E9E1F7", color="#A88FD9", fontcolor="#45296B"];
%     ContractS -> AskS [penwidth=1.6, color="#8592A3"];
%   }
% 
%   subgraph cluster_redundant {
%     label     = "Redundant match (k_red > 1)";
%     labelloc  = "t";
%     fontname  = "Helvetica-Bold";
%     fontsize  = 10.5;
%     fontcolor = "#3A4A5C";
%     style     = "rounded,filled";
%     fillcolor = "#F7F9FB";
%     color     = "#C7D0DA";
%     penwidth  = 1.0;
%     margin    = 14;
% 
%     ContractR [label=<<b>contract κ</b>>,
%                fillcolor="#DFEDE0", color="#8FB79A", fontcolor="#2E5A3D"];
%     Ask1      [label=<<b>staked ask 1</b>>,
%                fillcolor="#E9E1F7", color="#A88FD9", fontcolor="#45296B"];
%     Ask2      [label=<<b>staked ask 2</b>>,
%                fillcolor="#E9E1F7", color="#A88FD9", fontcolor="#45296B"];
%     Ask3      [label=<<b>staked ask k_red</b>>,
%                fillcolor="#E9E1F7", color="#A88FD9", fontcolor="#45296B"];
%     ContractR -> Ask1 [label="Race", penwidth=1.6, color="#8592A3"];
%     ContractR -> Ask2 [label="Race", penwidth=1.6, color="#8592A3"];
%     ContractR -> Ask3 [label="Race", penwidth=1.6, color="#8592A3"];
%   }
% }
% ```
% label=fig:redundancy_racing
% caption=Redundancy as a decentralization safeguard.
% rendered_images:end
% render_images:begin
\begin{figure}[H]
  \includegraphics[width=\linewidth]{08_decentralized_extension.tex.figs/08_decentralized_extension.5.png}
  \caption{Redundancy as a decentralization safeguard.}
  \label{fig:redundancy_racing}
\end{figure}
% render_images:end

% ==============================================================================
\subsection{Comparison with the centralized design}
\label{sec:decentralized_comparison}

Table~\ref{tab:decentralized_comparison} summarizes the trade-off between
the centralized \NoesisMarket{} of Section~\ref{sec:noesis_market} and the
decentralized variant sketched above.

\begin{samepage}
  \begin{table}[H]
    \centering
    \begin{tabular}{l|c|c|}
                                                & \emph{Centralized \NoesisMarket} & \emph{Decentralized \NoesisMarket} \\
      \hline
      Clearing requires a trusted operator      & \nbad     & \ngood    \\
      \hline
      Bidder identity and volume kept private   & \nbad     & \ngood    \\
      \hline
      Non-performance penalty is automatic      & \npartial & \ngood    \\
      \hline
      Settlement latency and cost               & \ngood    & \npartial \\
      \hline
      Capability quality independently verified & \nbad     & \nbad     \\
      \hline
    \end{tabular}
    \caption{Trade-offs between the centralized \NoesisMarket{} of
      Section~\ref{sec:noesis_market} and the decentralized, blind-bid
      variant sketched in this section.}
    \label{tab:decentralized_comparison}
  \end{table}
\end{samepage}

The last row is deliberately unchanged across both designs: staking
(Section~\ref{sec:stake_slashing}) makes non-fulfillment costly, but it
does not, by itself, make the delivered capability $\hat c(x)$
(Section~\ref{sec:fulfillment_monitoring}) any more verifiable than it is
in the centralized design. This is the same gap
Section~\ref{sec:related_dex} identifies relative to Tulip: a token
transfer settles atomically on-chain, but an LLM response's quality is
estimated, not settled, whichever design clears the contract that produced
it. Section~\ref{sec:decentralized_open_questions} collects the open
questions specific to this design.
